\documentclass[11pt,a4paper]{article}
\usepackage{amsmath,amssymb}
\usepackage{listings}
\usepackage{hyperref}
\usepackage[margin=1in]{geometry}
\usepackage{fancyhdr}

\pagestyle{fancy}
\fancyhf{}
\fancyhead[L]{Module 12: Advanced Data Structures}
\fancyhead[R]{C Programming Course}
\fancyfoot[C]{\thepage}

\lstset{
  language=C,
  basicstyle=\ttfamily\small,
  keywordstyle=\bfseries,
  commentstyle=\itshape,
  numbers=left,
  numberstyle=\tiny,
  frame=single,
  breaklines=true,
  tabsize=4
}

\title{Module 12: Advanced Data Structures}
\author{C Programming Course}
\date{}

\begin{document}
\maketitle
\thispagestyle{fancy}

\section{Binary Trees and BSTs}
\begin{lstlisting}
#include <stdio.h>
#include <stdlib.h>

typedef struct TreeNode {
    int key;
    struct TreeNode *left;
    struct TreeNode *right;
} TreeNode;

TreeNode *insert(TreeNode *root, int key) {
    if (root == NULL) {
        TreeNode *n = malloc(sizeof(TreeNode));
        if (n == NULL) return NULL;
        n->key = key;
        n->left = n->right = NULL;
        return n;
    }
    if (key < root->key)
        root->left = insert(root->left, key);
    else if (key > root->key)
        root->right = insert(root->right, key);
    return root;
}

void inorder(const TreeNode *root) {
    if (root == NULL) return;
    inorder(root->left);
    printf("%d ", root->key);
    inorder(root->right);
}
\end{lstlisting}

\section{AVL Trees and Red-Black Trees}
\begin{itemize}
  \item \textbf{AVL Tree} -- self-balancing BST; heights of left and right
    subtrees differ by at most 1. Uses rotations to rebalance.
  \item \textbf{Red-Black Tree} -- each node is colored red or black.
    Guarantees $O(\log n)$ operations with fewer rotations than AVL.
\end{itemize}

\section{Heaps and Priority Queues}
A min-heap ensures the parent is smaller than its children.
Implemented efficiently with an array:
\begin{itemize}
  \item Parent of node $i$: $(i - 1) / 2$
  \item Left child: $2i + 1$, Right child: $2i + 2$
\end{itemize}

\section{Graphs}
Represented as adjacency matrices or adjacency lists.
\begin{lstlisting}
#define MAX_V 100

typedef struct {
    int adj[MAX_V][MAX_V];
    int num_vertices;
} Graph;

void add_edge(Graph *g, int u, int v) {
    g->adj[u][v] = 1;
    g->adj[v][u] = 1;
}
\end{lstlisting}

Traversal algorithms: BFS (breadth-first) and DFS (depth-first).

\section{Tries}
A trie stores strings character by character. Each node has up to 26 children
(for lowercase English). Useful for prefix matching and autocomplete.

\end{document}
